\section{Introduction}\label{sec:intro}

Scene graph generation is evaluated on a benchmark whose labels are unusually predictable
from something every model sees. Given the ordered object pair (\textsf{man},
\textsf{surfboard}), the predicate on Visual Genome~\cite{krishna2017visual} is very often
\textsf{riding}, and a lookup table of training frequencies is famously
competitive~\cite{zellers2018neural}. Recall@$K$ rewards models that echo this prior, so the
field now leads with mean recall, which averages recall over predicates and so rewards getting
the rare ones~\cite{chen2019kern,tang2019vctree}. Debiasing methods are now judged by how far they raise it.
Total Direct Effect (TDE)~\cite{tang2020unbiased}, among the most cited of them, raises
mean recall@50 of a MOTIFS model from $14.6$ to $24.8$ on the canonical VG150 split, in our
evaluation of the released checkpoint with the authors' own code.

What did those ten points buy? There are two ways to raise mean recall, and the metric
cannot tell them apart. A model can move probability towards rare predicates \emph{within
each subject--object pair} --- say \textsf{carrying} more often for (\textsf{man},
\textsf{surfboard}) --- without becoming any better at deciding \emph{which} images show a
man carrying rather than riding. Or it can learn to make that decision: put more probability
on the correct predicate for the individual image, beyond what the object pair implies. The
first is a change in the model's group-level predictions; only the second is better
recognition. Every existing remedy --- head/body/tail-stratified recall~\cite{li2021bgnn}, independent
mean recall~\cite{li2022rethinking},
zero-shot recall~\cite{lu2016visual}, rarity-stratified accuracy~\cite{kervadec2021roses}
--- slices one model's performance differently. None attributes the gain \emph{between} two
models to one route or the other.

\paragraph{An exact split.} We separate the two routes with one operation
(\cref{fig:concept}). Take two frozen models and a test set grouped by a discrete feature
both models observe, here the subject--object pair. Score each model's prediction for one
relation against the label of a \emph{different} relation in the same group. This
relabelling preserves each group's predicate frequencies exactly, so any part of the score
gain earned at the level of the group survives it; it destroys the correspondence between an
image and its own label, so any part earned by getting individual cases right does not. The
observed gain therefore splits exactly into a \emph{group-level} part $\hdP$ and a
\emph{case-level} remainder $\hdR$, which is a within-group covariance between the change in
the models' predictions and the label (\cref{thm:population}). The split holds in the observed
sample for any proper score, needs no fitted reference model, held-out fold or tuning
parameter, costs $O(NK)$ on cached predictions, and comes with image-clustered intervals.

\paragraph{What it finds.} On the two models of TDE's own released checkpoint, evaluated
with the authors' code, and with each model's confidence matched on held-out images, the
proper-score change is almost entirely group-level: the case-level part is under a tenth
of it under either the quadratic or the log score, and under the quadratic score it is
$-0.00006$ with an interval straddling zero (\cref{sec:sggaudit}). The proper-score change behind the ten points is a change in which predicates the model
favours within each pair, far more than in how well it tells that pair's relations apart. Nor
is that implied by how TDE is built: in the released model its counterfactual subtraction
cancels the image-specific visual logits as well as the frequency prior (\cref{sec:sggaudit}),
so it discards case-level evidence along with the group-level prior. The converse appears too. A model reading image crops through a frozen
CLIP encoder~\cite{radford2021learning}, which every aggregate measure ranks last, gains
$+0.00467$ of case-level covariance over the geometry-based model ranked above it, with
calibration matched and an interval excluding zero (\cref{sec:vgvisual}): its deficit is
entirely group-level. That gain does not make it the better ranker, though, and the two
measures' disagreement is itself part of what \cref{sec:vgvisual} reports.

\paragraph{Contributions.}
\begin{itemize}
\setlength\itemsep{1pt}
\item An audit of a published debiasing method's released checkpoints, finding its ten-point
  mean-recall gain backed by a proper-score change that is almost entirely group-level --- a
  case-level part indistinguishable from zero under the quadratic score and under a tenth of
  the total under the log score --- and a reversal in which the model ranked last is the best
  discriminator among relations sharing an object pair.
\item An exact, model-free decomposition of the score gain between two frozen models into
  group-level and case-level parts, valid for every Bregman score, computed in one pass over
  cached predictions, with cluster-robust inference.
\item Validation that the case-level part is exactly zero when it should be, covers at the
  nominal rate under image clustering, and separates calibration, prior and shortcut
  effects --- together with a protocol for the one confound it does not separate.
\end{itemize}

\paragraph{What is not claimed.} The split is conditional on the declared grouping: $\hdR$
measures case-level discrimination \emph{among relations sharing an object pair}, and
credits any signal varying inside a pair, including a shortcut nobody recorded. It is not
invariant to recalibration --- rescaling a model's confidence moves score between the two
parts while changing no ranking --- which is why the comparisons whose models differ visibly in
confidence, TDE's and the CLIP model's, are read confidence-matched.
The split is of a proper score that weights every relation equally, whereas mean recall
weights every predicate equally, so the audit speaks to the relation-weighted change; a
case-level gain confined to rare predicates could be hidden by losses on common ones. And the
audit's headline holds at the configuration we argue for, not at every configuration; \cref{tab:sggaudit} reports all five. These limits are taken up
in \cref{sec:discussion}.
